% Native vector figure. Requires tikz and libraries arrows.meta,calc.
% Use inside a figure environment: \centering% Native vector figure. Requires tikz and libraries arrows.meta,calc.
% Use inside a figure environment: \centering% Native vector figure. Requires tikz and libraries arrows.meta,calc.
% Use inside a figure environment: \centering% Native vector figure. Requires tikz and libraries arrows.meta,calc.
% Use inside a figure environment: \centering\input{figures/modeling_framework}
\begingroup
\definecolor{MMFInk}{HTML}{273B48}
\definecolor{MMFLine}{HTML}{566570}
\definecolor{MMFPale}{HTML}{F1F4F5}
\begin{tikzpicture}[
  x=1cm,y=1cm,
  font=\fontsize{9}{12}\selectfont,
  ink/.style={text=MMFInk},
  block/.style={ink,draw=MMFLine,line width=.55pt,fill=white,
    align=center,inner xsep=3pt,inner ysep=4pt,text width=3.23cm,
    minimum width=3.5cm},
  foundation/.style={block,fill=MMFPale,minimum height=1.5cm},
  question/.style={block,minimum height=2cm},
  outcome/.style={block,minimum height=1.5cm},
  flow/.style={draw=MMFLine,line width=.6pt,-{Stealth[length=1.8mm,width=1.3mm]}},
  bus/.style={draw=MMFLine,line width=.6pt}
]
  % Common units, information boundary, physical model and settlement.
  \node[foundation] (input) at (-5.625,0)
    {\textbf{附件输入}\\负荷、光伏功率与电价\\已发布的日内预报};
  \node[foundation] (units) at (-1.875,0)
    {\textbf{量纲与信息边界}\\10分钟功率转电量\\仅用决策时可得信息};
  \node[foundation,text width=6.98cm,minimum width=7.25cm] (mother) at (3.75,0)
    {\textbf{统一母模型}\\供给$\geq$净需求；储能容量、功率与效率约束\\计划量全额结算；调整与紧急补购单独计费};
  \draw[flow] (input.east) -- (units.west);
  \draw[flow] (units.east) -- (mother.west);

  % The four question settings share the mother model.
  \node[question] (q1) at (-5.625,-2.55)
    {\textbf{问题一：确定性调度}\\已知负荷、光伏与电价\\购电与储能联合LP\\确定性最优调度基准};
  \node[question] (q2) at (-1.875,-2.55)
    {\textbf{问题二：日前采购}\\整日历史残差构造场景\\情景LP制定共同计划\\计划后按实际供需执行};
  \node[question] (q3) at (1.875,-2.55)
    {\textbf{问题三：滚动采购}\\固定规则回放候选\\0、6、12、18时筛选\\更新剩余购电合同};
  \node[question] (q4) at (5.625,-2.55)
    {\textbf{问题四：动态电价}\\同日配对价格与净负荷\\日前版沿用母模型\\滚动版沿用候选筛选};
  \coordinate (split) at (3.75,-1.15);
  \draw[bus] (mother.south) -- (split);
  \draw[bus] (-5.625,-1.15) -- (5.625,-1.15);
  \foreach \n in {q1,q2,q3,q4}
    \draw[flow] (\n.north |- split) -- (\n.north);

  % Q1 is a joint deterministic LP, not a run of the frozen rule executor.
  % The closed-loop execution and screening evaluation apply to Q2--Q4.
  \coordinate (collect) at (0,-4.02);
  \foreach \n in {q2,q3,q4}
    \draw[bus] (\n.south) -- (\n.south |- collect);
  \draw[bus] (-5.625,-4.02) -- (5.625,-4.02);
  \node[ink,align=center,font=\fontsize{9}{12}\selectfont] at (1.3,-4.5)
    {\textbf{问题二至四：闭环执行与开发期选型}};

  \node[outcome] (execute) at (-5.625,-5.65)
    {\textbf{固定实时执行器}\\每10分钟依实际供需\\充放电与紧急补购};
  \node[outcome] (evaluate) at (-1.875,-5.65)
    {\textbf{连续状态与实际账单}\\储电量跨日连续推进\\计划、调整及紧急费};
  \node[outcome] (ablate) at (1.875,-5.65)
    {\textbf{开发期配置对照}\\滚动：四种筛选时点\\费用与紧急补购权衡};
  \node[outcome,fill=MMFPale] (final) at (5.625,-5.65)
    {\textbf{最终分问配置}\\日前两组保留母模型\\滚动两组全时点筛选};
  \draw[flow] (-5.625,-4.02) -- (execute.north);
  \draw[flow] (execute.east) -- (evaluate.west);
  \draw[flow] (evaluate.east) -- (ablate.west);
  \draw[flow] (ablate.east) -- (final.west);
\end{tikzpicture}
\endgroup

\begingroup
\definecolor{MMFInk}{HTML}{273B48}
\definecolor{MMFLine}{HTML}{566570}
\definecolor{MMFPale}{HTML}{F1F4F5}
\begin{tikzpicture}[
  x=1cm,y=1cm,
  font=\fontsize{9}{12}\selectfont,
  ink/.style={text=MMFInk},
  block/.style={ink,draw=MMFLine,line width=.55pt,fill=white,
    align=center,inner xsep=3pt,inner ysep=4pt,text width=3.23cm,
    minimum width=3.5cm},
  foundation/.style={block,fill=MMFPale,minimum height=1.5cm},
  question/.style={block,minimum height=2cm},
  outcome/.style={block,minimum height=1.5cm},
  flow/.style={draw=MMFLine,line width=.6pt,-{Stealth[length=1.8mm,width=1.3mm]}},
  bus/.style={draw=MMFLine,line width=.6pt}
]
  % Common units, information boundary, physical model and settlement.
  \node[foundation] (input) at (-5.625,0)
    {\textbf{附件输入}\\负荷、光伏功率与电价\\已发布的日内预报};
  \node[foundation] (units) at (-1.875,0)
    {\textbf{量纲与信息边界}\\10分钟功率转电量\\仅用决策时可得信息};
  \node[foundation,text width=6.98cm,minimum width=7.25cm] (mother) at (3.75,0)
    {\textbf{统一母模型}\\供给$\geq$净需求；储能容量、功率与效率约束\\计划量全额结算；调整与紧急补购单独计费};
  \draw[flow] (input.east) -- (units.west);
  \draw[flow] (units.east) -- (mother.west);

  % The four question settings share the mother model.
  \node[question] (q1) at (-5.625,-2.55)
    {\textbf{问题一：确定性调度}\\已知负荷、光伏与电价\\购电与储能联合LP\\确定性最优调度基准};
  \node[question] (q2) at (-1.875,-2.55)
    {\textbf{问题二：日前采购}\\整日历史残差构造场景\\情景LP制定共同计划\\计划后按实际供需执行};
  \node[question] (q3) at (1.875,-2.55)
    {\textbf{问题三：滚动采购}\\固定规则回放候选\\0、6、12、18时筛选\\更新剩余购电合同};
  \node[question] (q4) at (5.625,-2.55)
    {\textbf{问题四：动态电价}\\同日配对价格与净负荷\\日前版沿用母模型\\滚动版沿用候选筛选};
  \coordinate (split) at (3.75,-1.15);
  \draw[bus] (mother.south) -- (split);
  \draw[bus] (-5.625,-1.15) -- (5.625,-1.15);
  \foreach \n in {q1,q2,q3,q4}
    \draw[flow] (\n.north |- split) -- (\n.north);

  % Q1 is a joint deterministic LP, not a run of the frozen rule executor.
  % The closed-loop execution and screening evaluation apply to Q2--Q4.
  \coordinate (collect) at (0,-4.02);
  \foreach \n in {q2,q3,q4}
    \draw[bus] (\n.south) -- (\n.south |- collect);
  \draw[bus] (-5.625,-4.02) -- (5.625,-4.02);
  \node[ink,align=center,font=\fontsize{9}{12}\selectfont] at (1.3,-4.5)
    {\textbf{问题二至四：闭环执行与开发期选型}};

  \node[outcome] (execute) at (-5.625,-5.65)
    {\textbf{固定实时执行器}\\每10分钟依实际供需\\充放电与紧急补购};
  \node[outcome] (evaluate) at (-1.875,-5.65)
    {\textbf{连续状态与实际账单}\\储电量跨日连续推进\\计划、调整及紧急费};
  \node[outcome] (ablate) at (1.875,-5.65)
    {\textbf{开发期配置对照}\\滚动：四种筛选时点\\费用与紧急补购权衡};
  \node[outcome,fill=MMFPale] (final) at (5.625,-5.65)
    {\textbf{最终分问配置}\\日前两组保留母模型\\滚动两组全时点筛选};
  \draw[flow] (-5.625,-4.02) -- (execute.north);
  \draw[flow] (execute.east) -- (evaluate.west);
  \draw[flow] (evaluate.east) -- (ablate.west);
  \draw[flow] (ablate.east) -- (final.west);
\end{tikzpicture}
\endgroup

\begingroup
\definecolor{MMFInk}{HTML}{273B48}
\definecolor{MMFLine}{HTML}{566570}
\definecolor{MMFPale}{HTML}{F1F4F5}
\begin{tikzpicture}[
  x=1cm,y=1cm,
  font=\fontsize{9}{12}\selectfont,
  ink/.style={text=MMFInk},
  block/.style={ink,draw=MMFLine,line width=.55pt,fill=white,
    align=center,inner xsep=3pt,inner ysep=4pt,text width=3.23cm,
    minimum width=3.5cm},
  foundation/.style={block,fill=MMFPale,minimum height=1.5cm},
  question/.style={block,minimum height=2cm},
  outcome/.style={block,minimum height=1.5cm},
  flow/.style={draw=MMFLine,line width=.6pt,-{Stealth[length=1.8mm,width=1.3mm]}},
  bus/.style={draw=MMFLine,line width=.6pt}
]
  % Common units, information boundary, physical model and settlement.
  \node[foundation] (input) at (-5.625,0)
    {\textbf{附件输入}\\负荷、光伏功率与电价\\已发布的日内预报};
  \node[foundation] (units) at (-1.875,0)
    {\textbf{量纲与信息边界}\\10分钟功率转电量\\仅用决策时可得信息};
  \node[foundation,text width=6.98cm,minimum width=7.25cm] (mother) at (3.75,0)
    {\textbf{统一母模型}\\供给$\geq$净需求；储能容量、功率与效率约束\\计划量全额结算；调整与紧急补购单独计费};
  \draw[flow] (input.east) -- (units.west);
  \draw[flow] (units.east) -- (mother.west);

  % The four question settings share the mother model.
  \node[question] (q1) at (-5.625,-2.55)
    {\textbf{问题一：确定性调度}\\已知负荷、光伏与电价\\购电与储能联合LP\\确定性最优调度基准};
  \node[question] (q2) at (-1.875,-2.55)
    {\textbf{问题二：日前采购}\\整日历史残差构造场景\\情景LP制定共同计划\\计划后按实际供需执行};
  \node[question] (q3) at (1.875,-2.55)
    {\textbf{问题三：滚动采购}\\固定规则回放候选\\0、6、12、18时筛选\\更新剩余购电合同};
  \node[question] (q4) at (5.625,-2.55)
    {\textbf{问题四：动态电价}\\同日配对价格与净负荷\\日前版沿用母模型\\滚动版沿用候选筛选};
  \coordinate (split) at (3.75,-1.15);
  \draw[bus] (mother.south) -- (split);
  \draw[bus] (-5.625,-1.15) -- (5.625,-1.15);
  \foreach \n in {q1,q2,q3,q4}
    \draw[flow] (\n.north |- split) -- (\n.north);

  % Q1 is a joint deterministic LP, not a run of the frozen rule executor.
  % The closed-loop execution and screening evaluation apply to Q2--Q4.
  \coordinate (collect) at (0,-4.02);
  \foreach \n in {q2,q3,q4}
    \draw[bus] (\n.south) -- (\n.south |- collect);
  \draw[bus] (-5.625,-4.02) -- (5.625,-4.02);
  \node[ink,align=center,font=\fontsize{9}{12}\selectfont] at (1.3,-4.5)
    {\textbf{问题二至四：闭环执行与开发期选型}};

  \node[outcome] (execute) at (-5.625,-5.65)
    {\textbf{固定实时执行器}\\每10分钟依实际供需\\充放电与紧急补购};
  \node[outcome] (evaluate) at (-1.875,-5.65)
    {\textbf{连续状态与实际账单}\\储电量跨日连续推进\\计划、调整及紧急费};
  \node[outcome] (ablate) at (1.875,-5.65)
    {\textbf{开发期配置对照}\\滚动：四种筛选时点\\费用与紧急补购权衡};
  \node[outcome,fill=MMFPale] (final) at (5.625,-5.65)
    {\textbf{最终分问配置}\\日前两组保留母模型\\滚动两组全时点筛选};
  \draw[flow] (-5.625,-4.02) -- (execute.north);
  \draw[flow] (execute.east) -- (evaluate.west);
  \draw[flow] (evaluate.east) -- (ablate.west);
  \draw[flow] (ablate.east) -- (final.west);
\end{tikzpicture}
\endgroup

\begingroup
\definecolor{MMFInk}{HTML}{273B48}
\definecolor{MMFLine}{HTML}{566570}
\definecolor{MMFPale}{HTML}{F1F4F5}
\begin{tikzpicture}[
  x=1cm,y=1cm,
  font=\fontsize{9}{12}\selectfont,
  ink/.style={text=MMFInk},
  block/.style={ink,draw=MMFLine,line width=.55pt,fill=white,
    align=center,inner xsep=3pt,inner ysep=4pt,text width=3.23cm,
    minimum width=3.5cm},
  foundation/.style={block,fill=MMFPale,minimum height=1.5cm},
  question/.style={block,minimum height=2cm},
  outcome/.style={block,minimum height=1.5cm},
  flow/.style={draw=MMFLine,line width=.6pt,-{Stealth[length=1.8mm,width=1.3mm]}},
  bus/.style={draw=MMFLine,line width=.6pt}
]
  % Common units, information boundary, physical model and settlement.
  \node[foundation] (input) at (-5.625,0)
    {\textbf{附件输入}\\负荷、光伏功率与电价\\已发布的日内预报};
  \node[foundation] (units) at (-1.875,0)
    {\textbf{量纲与信息边界}\\10分钟功率转电量\\仅用决策时可得信息};
  \node[foundation,text width=6.98cm,minimum width=7.25cm] (mother) at (3.75,0)
    {\textbf{统一母模型}\\供给$\geq$净需求；储能容量、功率与效率约束\\计划量全额结算；调整与紧急补购单独计费};
  \draw[flow] (input.east) -- (units.west);
  \draw[flow] (units.east) -- (mother.west);

  % The four question settings share the mother model.
  \node[question] (q1) at (-5.625,-2.55)
    {\textbf{问题一：确定性调度}\\已知负荷、光伏与电价\\购电与储能联合LP\\确定性最优调度基准};
  \node[question] (q2) at (-1.875,-2.55)
    {\textbf{问题二：日前采购}\\整日历史残差构造场景\\情景LP制定共同计划\\计划后按实际供需执行};
  \node[question] (q3) at (1.875,-2.55)
    {\textbf{问题三：滚动采购}\\固定规则回放候选\\0、6、12、18时筛选\\更新剩余购电合同};
  \node[question] (q4) at (5.625,-2.55)
    {\textbf{问题四：动态电价}\\同日配对价格与净负荷\\日前版沿用母模型\\滚动版沿用候选筛选};
  \coordinate (split) at (3.75,-1.15);
  \draw[bus] (mother.south) -- (split);
  \draw[bus] (-5.625,-1.15) -- (5.625,-1.15);
  \foreach \n in {q1,q2,q3,q4}
    \draw[flow] (\n.north |- split) -- (\n.north);

  % Q1 is a joint deterministic LP, not a run of the frozen rule executor.
  % The closed-loop execution and screening evaluation apply to Q2--Q4.
  \coordinate (collect) at (0,-4.02);
  \foreach \n in {q2,q3,q4}
    \draw[bus] (\n.south) -- (\n.south |- collect);
  \draw[bus] (-5.625,-4.02) -- (5.625,-4.02);
  \node[ink,align=center,font=\fontsize{9}{12}\selectfont] at (1.3,-4.5)
    {\textbf{问题二至四：闭环执行与开发期选型}};

  \node[outcome] (execute) at (-5.625,-5.65)
    {\textbf{固定实时执行器}\\每10分钟依实际供需\\充放电与紧急补购};
  \node[outcome] (evaluate) at (-1.875,-5.65)
    {\textbf{连续状态与实际账单}\\储电量跨日连续推进\\计划、调整及紧急费};
  \node[outcome] (ablate) at (1.875,-5.65)
    {\textbf{开发期配置对照}\\滚动：四种筛选时点\\费用与紧急补购权衡};
  \node[outcome,fill=MMFPale] (final) at (5.625,-5.65)
    {\textbf{最终分问配置}\\日前两组保留母模型\\滚动两组全时点筛选};
  \draw[flow] (-5.625,-4.02) -- (execute.north);
  \draw[flow] (execute.east) -- (evaluate.west);
  \draw[flow] (evaluate.east) -- (ablate.west);
  \draw[flow] (ablate.east) -- (final.west);
\end{tikzpicture}
\endgroup
